% ===================================================================
%  بخش ۱ — مقدمه
%  منبع: متن موجود مقاله (`first artcel text.txt`)
%  + دو پاراگراف پایانی نو از `draft/section1_intro_addition.md` (کار ۲)
% ===================================================================
\section{مقدمه}
\label{sec:intro}

دستکاری نرخ گسیل خودبه‌خودی اتم‌ها، مولکول‌های فلورسنت و نقاط کوانتومی از
طریق اصلاح چگالی موضعی حالات فوتونی (LDOS)، موضوعی محوری در الکترودینامیک
کوانتومی کاواکی (cQED) و نانوفوتونیک مدرن است. بر اساس اثر پورسل
\cite{purcell1946spontaneous}، نرخ گسیل خودبه‌خودی یک گسیل‌گر می‌تواند با
قرارگیری در داخل یا مجاورت یک تشدیدگر اپتیکی به شدت افزایش یابد. در حالی
که کاواک‌های دی‌الکتریک سنتی به دلیل حد پراش نور دارای حجم مِدی نسبتاً
بزرگی هستند، نانوساختارهای فلزی با تکیه بر تشدید پلاسمون‌های سطحی موضعی
(LSPR) قادرند نور را در مقیاس‌های بسیار کوچک‌تر از طول‌موج فشرده کرده و
حجم مدی فوق‌العاده پایینی ($V_{\text{mode}} \ll \lambda^3$) فراهم آورند
\cite[pp.~34--37]{maier2007plasmonics}. این وابستگی به محیط تنها به
ساختار فوتونی پیرامون محدود نمی‌شود؛ در گسیل‌گرهای حالت جامد، جزئیات
محیط مادی نیز بر مشخصه‌های گسیل اثر می‌گذارد — نمونه‌ی آن کاهش کیفیت
درهم‌تنیدگی قطبش فوتون‌های حاصل از فروافت آبشاری دواکسیتون در یک نقطه‌ی
کوانتومی خودتجمع‌یافته، در اثر لایه‌ی مرطوب‌کننده و افزایش دما است
\cite{bagheriharouni2010wetting}.

در میان هندسه‌های مختلف نانوذرات فلزی، نانومیله‌های طلا (Gold Nanorods) به
دلیل قابلیت تنظیم طول‌موج تشدید طولی در بازه مرئی و فروسرخ نزدیک و
همچنین ایجاد میدان‌های به شدت متمرکز در نوک‌های دارای انحنا (اثر رعدوبرق
یا Lightning-rod effect)، جایگاه ویژه‌ای در طراحی نانوآنتن‌های نوری دارند
\cite{chen2009subwavelength}. با این حال، قرارگیری گسیل‌گر در فاصله بسیار
نزدیک از سطح فلز چالش رقابت میان تقویت تابشی و اتلافات اهمی ناشی از جذب
درون‌باندی فلز (پدیده خاموشی یا Quenching) را به همراه دارد
\cite{anger2006enhancement, koenderink2017single}. از این رو، مدل‌سازی
دقیق سهم کانال‌های تابشی و غیرتابشی نقشی کلیدی در بهینه‌سازی عملکرد این
سیستم‌ها ایفا می‌کند. در این کار، با پیاده‌سازی یک مدل FDTD متقارن‌محوری دقیق،
پاسخ الکترومغناطیسی، تقویت فاکتور پورسل، راندمان تابشی و نقشه میدان نزدیک
نانومیله طلا به تفکیک استخراج و تحلیل شده است.

% -------------------------------------------------------------------
% >>> کار ۲: دو پاراگراف افزودنی — draft/section1_intro_addition.md
% -------------------------------------------------------------------

زبانی که امروز برای توصیف گسیل خودبه‌خودی به کار می‌رود، از فیزیک اتمی به
ارث رسیده است. در این زبان، نرخ گذار میان دو تراز با ضرایب اینشتین —
ضریب گسیل خودبه‌خودی $A_{21}$ و ضرایب گسیل و جذب القایی — بیان می‌شود و
ترازها با برچسب‌های $s$، $p$ و $d$ شناسانده می‌گردند. این چارچوب بر تقارن
کروی پتانسیل کولنی بنا شده است: هارمونیک‌های کروی
$Y_{lm}(\theta,\varphi)$ و عددهای $l$ و $m$، نتیجه‌ی جداسازی معادله‌ی موج
در دستگاه مختصات کروی‌اند و اعتبارشان به همین تقارن گره خورده است
\cite{griffiths2018introduction, jackson1999classical}. همین زبان، تقریباً
بی‌تغییر، به نانوفوتونیک منتقل شده است؛ مدهای پلاسمونی نانوذره‌ی کروی با
همان برچسب‌های $l$ و $m$ رده‌بندی می‌شوند و مدل هیبریداسیون پلاسمونی — که
پاسخ نانوساختارهای مرکب را به‌خوبی پیش‌بینی می‌کند — صریحاً از قیاس
نظریه‌ی اوربیتال مولکولی بهره می‌گیرد \cite{prodan2003hybridization}.

انتقال این چارچوب به هندسه‌ی نانومیله دو پرسش را پیش می‌کشد. نخست،
نانومیله تقارن کروی ندارد و محور طولی آن یک محور ممتاز است؛ نشانه‌ی
تجربی این تفاوت آشکار است، چرا که نانوذره‌ی کروی \textbf{یک} تشدید
پلاسمونی دارد و نانومیله به \textbf{دو} تشدید مجزای طولی و عرضی با پهنای
خط و انرژی متفاوت شکافته می‌شود \cite{sonnichsen2002drastic}. افزون بر
این، جایگاه ویژه‌ی کره در ادبیات موضوع بیش از هر چیز ریشه‌ای ریاضی دارد:
کره از معدود هندسه‌هایی است که مسئله‌ی پراکندگی در آن حل تحلیلی بسته
دارد. حل‌شدنی بودن، خاصیت دستگاه مختصات است و نه خاصیتی از فیزیک مسئله.
دوم، در فاصله‌های نانومتری از سطح فلز، کانال غیرتابشی فعال می‌شود و بازده
تابشی به فاصله و طول‌موج وابسته می‌گردد؛ اندازه‌گیری‌های تک‌مولکولی و
محاسبات نظری هر دو نشان داده‌اند که تقویت و خاموشی فلورسانس رفتار یکسانی
ندارند و آهنگ‌های تابشی و غیرتابشی از آغاز باید جداگانه فرمول‌بندی شوند
\cite{anger2006enhancement, bharadwaj2007spectral, carminati2006radiative}.
بر این پایه، در این پژوهش علاوه بر کمی‌سازی تقویت پورسل و راندمان تابشی
نانومیله‌ی طلا، نشان داده می‌شود که همین نتایج عددی مستقیماً دامنه‌ی
اعتبار توصیف‌های همسانگرد و تک‌پارامتری را آشکار می‌کنند. تأکید می‌شود که
هدف، ابطال این توصیف‌ها نیست بلکه تعیین مرزی است که هندسه‌ی ناهمسانگرد بر
آن‌ها تحمیل می‌کند.
